% Part: second-order-logic
% Chapter: syntax-and-semantics
% Section: terms-formulas

\documentclass[../../../include/open-logic-section]{subfiles}

\begin{document}

\olfileid[fr]{sol}{syn}{frm}

\olsection{Termes et \printtoken{p}{formula}}

Comme en logique du premier ordre, les expressions
de la logique du second ordre sont construites à
partir d'un vocabulaire de base comprenant des
\emph{!!{variable}s}, des \emph{!!{constant}s},
des \emph{!!{predicate}s} et parfois des
\emph{!!{function}s}. Avec les connecteurs
logiques, les quantificateurs et les signes de
ponctuation tels que les parenthèses et les
virgules, ces symboles permettent de former des
\emph{termes} et des \emph{!!{formula}s}.
La différence est que la logique du second ordre
possède, outre les variables d'objet, des
variables de relation et de fonction, sur lesquelles
elle permet également de quantifier. Ses symboles
logiques sont donc ceux du premier ordre, auxquels
s'ajoutent :

\begin{enumerate}
\item Pour chaque arité~$n$, un ensemble
  dénombrable de !!{variable}s de relation du
  second ordre : $\Obj{V^n_0}$, $\Obj{V^n_1}$,
  $\Obj{V^n_2}$, \dots
\item Pour chaque arité~$n$, un ensemble dénombrable de
  !!{variable}s de fonction du second ordre : $\Obj{u^n_0}$,
  $\Obj{u^n_1}$, $\Obj{u^n_2}$, \dots
\end{enumerate}

De même que $x$, $y$, $z$ sont des métavariables
pour les variables du premier ordre~$\Obj{v_i}$,
nous emploierons $X$, $Y$, $Z$, etc., comme
métavariables pour les variables de relation
$\Obj{V^n_i}$, et $u$, $v$, etc., comme
métavariables pour les variables de fonction
$\Obj{u^n_i}$.

\begin{explain}
On spécifie les symboles non logiques d'un langage
du second ordre comme ceux d'un langage du premier
ordre : en énumérant ses !!{constant}s, ses
!!{function}s et ses !!{predicate}s.

En logique du premier ordre, on inclut d'ordinaire
l'!!{identity}~$\eq$. Les symboles non logiques
d'un langage~$\Lang{L}$ y sont essentiels pour
exprimer des choses intéressantes. Il existe bien
des !!{sentence}s qui n'en utilisent aucun, mais
l'!!{identity} seule ne permet pas de dire
grand-chose. En logique du second ordre, la réserve
illimitée de variables de relation et de fonction
permet d'exprimer ce qu'exprimerait un langage du
premier ordre sans avoir à fixer un stock particulier
de symboles non logiques.\footnote{Les variables libres servent alors de paramètres d'interprétation pour les symboles remplacés.}
\end{explain}

\begin{defn}[Termes du second ordre]
L'ensemble des \emph{termes du second ordre} de
$\Lang L$, noté $\TrmSOL[L]$, s'obtient en
ajoutant à la \olref[fol][syn][frm]{defn:terms}
la clause suivante ; les clauses héritées s’appliquent aussi aux termes du second ordre :
\begin{enumerate}
\item Si $u$ est une variable de fonction à
  $n$ arguments et si $t_1$, \dots, $t_n$ sont
  des termes, alors $\Atom{u}{t_1,\ldots,t_n}$
  est un terme.
\end{enumerate}
\end{defn}

\begin{explain}
Un terme du second ordre ressemble donc à un
terme du premier ordre, sauf que, là où le terme
du premier ordre contient un
!!{function}~$\Obj{f^n_i}$, celui du second ordre
peut contenir à sa place une variable de
fonction~$\Obj{u^n_i}$.
\end{explain}

\begin{defn}[Formules du second ordre]
L'ensemble des \emph{!!{formula}s du second ordre}
$\FrmSOL[L]$ du langage~$\Lang L$ s'obtient en
ajoutant à la \olref[fol][syn][frm]{defn:formulas}
les clauses suivantes ; les clauses héritées s’appliquent aux termes du second ordre :
\begin{enumerate}
\item Si $X$ est une variable de prédicat à
  $n$ arguments et si $t_1$, \dots, $t_n$ sont
  des termes du second ordre de~$\Lang L$, alors
  $\Atom{X}{t_1,\ldots,t_n}$ est une
  !!{formula} atomique.

\tagitem{prvAll}{Si $!A$ est une !!{formula} et
  $u$ une variable de fonction, alors
  $\lforall[u][!A]$ est une !!{formula}.}{}

\tagitem{prvAll}{Si $!A$ est une !!{formula} et
  $X$ une variable de prédicat, alors
  $\lforall[X][!A]$ est une !!{formula}.}{}

\tagitem{prvEx}{Si $!A$ est une !!{formula} et
  $u$ une variable de fonction, alors
  $\lexists[u][!A]$ est une !!{formula}.}{}

\tagitem{prvEx}{Si $!A$ est une !!{formula} et
  $X$ une variable de prédicat, alors
  $\lexists[X][!A]$ est une !!{formula}.}{}
\end{enumerate}
\end{defn}

\end{document}
